% SOURCE docs/release/2026-10-29/source/docs/writing-repair/2026-10-29/src/6972ed269a657f6a.tex CHAR 0
\chapter{Operating a study without changing its scientific meaning}
\label{in:operations}

Operating an analysis preserves more than its final number. The argument begins with execution identity and recovery, then develops the publication, retention and capacity boundaries that make a result usable. Discovery, funding, observation and servicing examples show how those boundaries interact with causal uncertainty and monetary consequences. Their exact interface screens and mathematical constructions answer different questions from the fitted studies elsewhere in the book; the comparison keeps those evidential roles explicit.

\section{An execution identity is larger than a request name}

A durable analysis must bind the canonical request, scientific context, input
bytes, code, schema, numerical environment, seed and budget. A repeated job
name alone cannot establish that a result is still valid. The inspected runtime binds the request and a selected source/schema/environment fingerprint and rejects key collisions when those recorded identities change. The fuller identity is a design obligation; the developed recovery argument below identifies dependencies that the selected fingerprint does not independently record.~\citep{v12_in_runtime}

Artifacts use the specified Python JSON encoding, content and descriptor hashes, opaque identities and bounded local reads. Calling that encoding canonical describes this implementation's byte rule; it does not establish cross-language canonicalization. Regular-file and nonsymlink checks reduce accidental
path redirection. A content hash establishes byte identity under the relevant
assumption about the hash function; it is not an authenticated endorsement of
the input's scientific accuracy. A local operator can still supply a false
declaration whose bytes remain perfectly stable.

The scientific request contains no arbitrary shell command, Python import or
network endpoint. Fixed worker entry points execute the supported operations.
This makes the scope of the local tool explicit and allows a review grant to
refer to actual operations. The current controller's trusted caller object is
not a substitute for an institutional identity and approval service.

\section{Budgets, durable jobs and cancellation}

Scientific workers have wall-time and address-space limits. The memory bound
applies to the worker, not to all memory used by the resident SDK and
coordinator. A complete operating measurement must include both. Small control
or error envelopes may exceed an extremely small requested output limit to
explain why the request cannot be fulfilled. These are disclosed boundaries of
the present implementation.~\citep{v12_in_runtime}

Report jobs retain a hashed prefix of completed work identities. The local
retry rule caps worker starts at three and accounts for remaining time against the original local wall-clock budget. That rule is not a hard end-to-end deadline guarantee. Individual publication files are atomic, while artifact content, metadata and terminal job state remain separate transitions. Replay checks request, artifact and checkpoint identities before continuation. The recorded restart fixture kills a report worker after its
first checkpoint and completes all one hundred distinct work identities once.
This is evidence for that local workflow, not a distributed exactly-once
guarantee for arbitrary external side effects.

Cancelling an in-flight protocol request and cancelling a durable application
job are different operations. The inspected specification also recognizes
cancellation/completion races.~\citep{v12_mcp_legacy_cancellation} If publication
has already completed, a cancellation response must report that terminal state.
It must not imply that a displayed result or an external action was undone.
The present service has no operation that changes a bank account or sends a
customer communication.

\section{Fault evidence and what remains unmeasured}

Nine recorded fault cases cover replay with conflicting budgets, wall timeout,
output limits, worker memory, kill/restart with checkpoint, cooperative
cancellation, changed artifact bytes, changed scientific context and changed
checkpoints.~\citep{v12_in_runtime} The tests distinguish preserving a result from
merely restarting a process. An interrupted computation that silently uses new
inputs is a different analysis even if it eventually finishes.

These fixtures do not measure institution-scale availability, multi-host
recovery, adversarial tenant isolation or a service-level commitment. They also
do not show that streaming always improves elapsed time. The earlier complete
local workload comparisons retain small or absent gains; engineering choices
must be justified by the intended workload rather than by the existence of a
streaming interface.

Operational qualification should measure the complete request: artifact import,
scientific work, uncertainty, serialization, reporting and recovery. A benchmark
of one reducer omits fitting and coordination costs. The comparison needs
matched inputs, numerical tolerance and an explicit accounting of unsuccessful
requests, just as a scientific comparison needs a common target.

\section{Compatibility is a maintained scientific obligation}

A new backend release can change defaults, graph encodings, estimator targets,
uncertainty or failure behavior. A compatible JSON schema alone does not prove
scientific compatibility. The proposed adapter register should bind a version,
accepted inputs, method identity, preserved result types and matched adverse
fixtures. A material semantic change requires a reviewed successor rather than
relabeling an old qualification.

Transport migration likewise needs real client/server tests. The current legacy
MCP choice is useful for its tested environment. Supporting a modern protocol
or both eras requires a new compatibility argument, discovery behavior and
failure tests. It should not be advertised by extending a version string while
keeping the old handshake.

Ownership must also remain explicit. The causal library owns its result
contract and adapters. Numerical providers own their implementations and
scientific domains. Consumers own their target interpretation and displayed
decision. The bank owns actual data authority, assignment records, operational
constraints and acceptance. A report can connect those responsibilities; it
cannot replace their evidence.

\section{What a recovery identity must preserve}

Consider a daily deposit report that aggregates customer-level evidence into
a 90-day funding decision. A restarted process must answer the same question
using the same admitted observations. Replacing a data extract with a later
vintage, changing a sign convention, or upgrading a numerical routine can
alter the scientific result even when the file name and command remain
unchanged. This is why recovery needs an analysis identity rather than only
a process identifier.

Write the intended identity as
\[
 J=\operatorname{Hash}(Q,A,S,E,\rho,B),
\]
where $Q$ is the scientific request, $A$ the input identities, $S$ the schema
and method semantics, $E$ the relevant execution environment, $\rho$ the
randomization state where applicable, and $B$ the declared budget. This is a
proposed completeness condition. The inspected runtime fingerprints a selected
list of source files and schemas, Python, and four named package versions.
Its request digest omits the display request identifier and idempotency key
while retaining the other request fields. It rejects reuse of a key with a
changed request or selected execution identity
\citep{exec13_ops_source}.

The selected fingerprint is useful but does not describe every numerical or
platform dependency. It does not, for example, independently record the
complete transitive dependency graph, numerical-library build, operating
system and hardware state. A future numerical adapter must identify which
of those properties can affect its result, then bind and test them. Extending
the fingerprint is an engineering obligation; it does not establish that the
underlying causal assumptions are true. Conversely, irrelevant display
metadata need not invalidate a deterministic result if the scientific and
execution identities remain unchanged.

Three compatibility questions follow. Byte replay asks whether the same
admitted computation returns the same canonical bytes. Numerical compatibility
allows a declared tolerance or stochastic comparison under an unchanged
mathematical target. Scientific compatibility asks whether the population,
treatment, horizon, units, assumptions and inferential meaning remain the
same. A backend upgrade can satisfy a numerical tolerance while changing
which observations are dropped or which population is weighted. Such an
upgrade needs a semantic review even if its output schema is unchanged.

For the funding report, sensible adverse fixtures include a changed date
cutoff, a currency-unit conversion, a negative monetary map, a missing
assignment field, an unbounded uncertainty set and a default that drops
customers without complete histories. The comparison must preserve valid
refusals and ambiguous results. Substituting a point estimate for an
unbounded set is not a compatibility repair.

\section{Checkpoint invariants and the scope of idempotence}

Suppose a report combines deterministic contributions $s_i$ for a frozen
ordered list of work identities $(a_1,\ldots,a_m)$. A useful checkpoint is
\[
 C_k=(J,(a_1,\ldots,a_k),(s_1,\ldots,s_k),h_k),
\]
where $h_k$ binds the preceding contents. The required invariant is that the
completed identifiers form a prefix of the frozen request, each identifier
has one committed contribution, and the public result is a function of those
committed contributions. Recomputing an uncommitted contribution after a
crash may waste work; committing it twice changes the answer.

An exact project construction makes the distinction visible. Three intended
contributions are $100,-40,70$, so the total is $130$. A crash and retry can
produce the calculation trace $100,-40,-40,70$. Naively appending every
returned value gives $90$. A map from work identity to its immutable
contribution still gives $130$, provided a repeated identity cannot carry
changed contents. These are finite arithmetic results, not measurements of
a distributed service~\citep{exec13_ops_exact}.

The inspected local report worker verifies its completed prefix and digest,
skips already committed identities, and updates summaries and completed
identities together in the job state. Its preserved restart trial completed
100 distinct work identities, with 100 unique entries, after two worker
starts. The more recent preserved fault replay also reports the nine declared
cases passing; that replay was not a fresh MCP client/server trial
\citep{exec13_ops_source,exec13_ops_faults}. This evidence supports the named
local report path. It does not establish recovery for arbitrary estimators,
multi-host execution or transactions at a bank.

The atomicity boundary is also specific. The artifact store publishes its
content and metadata as individually atomic files, checks any already
present bytes, and can complete a repeated deterministic publication. The
report's result reference and terminal job state are then published in the
job record. These are several transitions, not one transaction spanning
arbitrary external systems. A crash can leave an unreferenced artifact even
when no incorrect public result is exposed. Retention and garbage collection
therefore need a separate rule that preserves artifacts referenced by admitted
jobs and evidence. Hash verification alone cannot establish that cleanup is
safe.

A customer communication illustrates the stronger problem. If an external
system sends an offer and the acknowledgement is lost, a retrying client may
not know whether the offer was sent. Recording the local calculation once
does not resolve that ambiguity. An eventual external-action adapter would
need a cooperating transaction or deduplication contract, reconciliation and
a named owner for unresolved outcomes. The present service performs no such
bank action. Its proposal should retain report generation as a distinct,
read-oriented capability rather than infer action safety from a successful
restart.

\section{Retries purchase attempts, not independent chances of success}

Let $K$ be a maximum number of attempts and $q$ the probability of a transient
failure on each attempt. If failures are independent and identically
distributed, the probability of obtaining a result is $1-q^K$ and the
expected number of attempts is
\[
 \mathbb E N_K=\sum_{j=0}^{K-1}q^j.
\]
For the stipulated values $q=1/5$ and $K=3$, these are $124/125=0.992$
and $31/25=1.24$. The numbers express an assumed failure law. Nine passing
fault fixtures do not estimate $q$, and a retry cap alone cannot establish
the 99.2\% success probability~\citep{exec13_ops_exact}.

Persistent failures change the calculation. Suppose one tenth of requests
cannot succeed under the supplied environment, while the remaining nine
tenths have independent per-attempt failure probability one tenth. Three
attempts then succeed with probability
\[
 \frac9{10}\left(1-\frac1{10^3}\right)=0.8991,
\]
and require $0.1(3)+0.9(1+0.1+0.01)=1.299$ attempts on average.
The marginal first-attempt failure probability is $0.19$. Treating that
marginal value as independent across retries would predict
$1-0.19^3=0.993141$, substantially overstating success. The retry has not
changed an unavailable dependency, an unsupported method or contradictory
input declarations. Operational classification should distinguish those
persistent causes from a recoverable interruption.

A bounded retry policy also retains its original time account. In the
inspected report path, remaining time is computed from the job creation
timestamp and original wall budget, and the worker-start count is capped at
three. Restarting does not intentionally grant three new full deadlines.
The implementation uses local wall-clock timestamps and operating-system
timers; this is not a distributed clock or a hard real-time guarantee.
Launch overhead, host suspension, clock adjustment and recovery of a failed
machine require a target-environment study before making a service-level
claim~\citep{exec13_ops_source}.

Statistical retries need an additional distinction. Replaying an interrupted
calculation with preserved random state is different from drawing new
randomness until a desired conclusion appears. If three independent valid
null $p$-values are generated and the smallest is selected, the probability
of at least one value below $0.05$ is
\[
 1-(1-0.05)^3=0.142625.
\]
This exact counterexample assumes independent uniform null $p$-values;
it does not allege that the runtime performs such selection. It shows why
the proposed adapter must separate mechanical recovery from a new statistical
replicate and preserve all requested analyses. The same issue arises when
timeouts depend on difficult data and only easy completed fits enter a
coverage report~\citep{exec13_ops_exact}.

\section{Resource budgets and cancellation at the right boundary}

The local worker sets an address-space limit and a wall-time timer. The
coordinator starts a fixed worker entry point, collects its output and
validates the canonical response against the requested output size. This
supports a bounded set of local operations. It does not impose one aggregate
memory bound on the worker, coordinator, SDK and all clients, and the output
check is applied after collection rather than as an incremental bound on an
arbitrary producer's stream~\citep{exec13_ops_source}.

For a proposed third-party numerical adapter, admission should therefore
include peak resident memory of the whole workflow, temporary disk use,
serialization and parent-process buffers. A limit that is safe for a short
report may be inappropriate for an estimator returning a large covariance
matrix. If a result exceeds the public response size, the design can publish
a separately bounded artifact with explicit custody and expose a small
reference. It must retain the scientific object's meaning and cannot silently
truncate interval components or omit unsuccessful fits to meet a size limit.

Cancellation is a race between a request and the computation's next state
transition. For pure report generation, a useful rule is that an accepted
cancellation recorded before the publication transition prevents a new public
result; cancellation after
publication returns the actual completed state. The inspected implementation
checks a cancellation flag around report work and before publication. That
cooperative behavior does not guarantee that every numerical kernel can be
interrupted at an arbitrary instruction. A later adapter may need worker
termination and a rule for discarding uncommitted artifacts.

A timeout should be recorded as an execution outcome. A scientifically
unidentified effect can instead be a completed result with an explicit
scientific limitation. Converting both into one generic failure statistic
obscures the remedy: the first may need resource or implementation work, while
the second may need a different target, new information or abstention. A
successful process exit likewise says nothing by itself about whether the
effect is precise enough to support a decision.

\section{Ownership and a reversible adoption path}

Operational ownership follows the claim being made. A method maintainer owns
the supported algorithm and numerical domain. The common-library owner owns
the request/result contract, method-version binding and admitted conversions.
The consumer owns the displayed target, units and business action. A later
bank study additionally needs a data owner for the actual observation and
assignment processes and a decision owner for losses and feasible actions.
These are proposed responsibilities, not names or approvals that have already
been supplied.

For example, a deposit report may use a valid backend estimate while a
consumer interprets the coefficient as currency rather than thousands of
currency units. The remedy belongs partly in the shared unit contract and
partly in the consuming application's acceptance examples. A central
numerical team cannot certify a silently altered presentation. Conversely,
an application owner cannot declare a confidence procedure valid merely
because the resulting interval is displayed correctly.

The proposal should make an adoption decision reversible at its natural
boundaries. Retain the incumbent expert workflow and its artifacts while a
new adapter is evaluated on matched public/generated requests. Publish both
methods' complete outputs with their own identities; compare discrepancies
before replacing the default. A schema migration should have a declared
mapping and refusal cases for information it cannot preserve. Where a new
method changes the estimand, the migration is a new scientific analysis,
not a formatting conversion. A rollback retains access to the previously
admitted version and evidence without rewriting its result as if produced
by the new version.

The same approach applies to workspace consumers. Read-only inspection of a
repository can reveal candidate objects and interfaces. It cannot establish
that its owner accepts the proposed contract, that its launch environment is
compatible, or that an integration saves work. A local mock consumer can test
the proposed request and response boundary now. Actual consumer execution,
deployment ownership and institutional operation remain distinct later
claims, with their own narrowly defined studies.

\section{The economic denominator is every attempted request}

Let $R$ denote whether a request returns a usable result and $D$ its incremental
decision value relative to the incumbent workflow, including the specified
fallback when $R=0$. Let $C$ include all request-specific compute, review,
retry and recovery costs. The per-request economic target is
\[
 \mathbb E[D-C],
\]
plus a consistent allocation of fixed construction and shared operating
costs. Reporting $\mathbb E[D\mid R=1]$ with a nominal successful-run cost
answers a different question. Reliability and economic comparisons require
both the success probability and what happens on the unsuccessful branch.

In an exact stipulated example, 70\% of requests are easy, complete and add
USD~10 of gross value. Another 15\% complete but add USD~$-30$. The remaining
15\% fail and revert to the baseline action, adding zero gross value but
requiring USD~20 of recovery work. Every attempt costs USD~2 in tools and
computation. The completed-only calculation reports
\[
 \frac{0.70(10)+0.15(-30)}{0.85}-2
 =\frac{16}{17}\simeq0.941
\]
USD per completed request. The all-attempt account instead gives
\[
 0.70(10)+0.15(-30)-2-0.15(20)=-2.5
\]
USD per attempted request. The apparent gain changes sign when the
unsuccessful work and correct denominator enter the account. These numbers
are generated planning assumptions, not workflow prices or measured bank
returns~\citep{exec13_ops_exact}.

Measurement should therefore retain attempted task identity, terminal
execution state, scientific disposition, usable-result status, fallback
action and disjoint cost categories. Repeated attempts are costs within a
request, not new independent banking opportunities. Correlated requests from
one institution, calendar or scientific task family also cannot be treated
as independent evidence merely because their request identifiers differ.
The analysis unit should follow the source of independent information.

A prototype can measure machine resources and trace whether required fields
survive recovery. It cannot price saved analyst time without observing a
matched workflow, nor infer business value from seconds saved. The next
inexpensive operational study should compare complete local workflows under
declared faults and budgets, preserve adverse results, and report source and
environment identities. Only a subsequent justified platform or bank study
can supply institution-specific availability, ownership, labor, and
decision-value evidence. Their absence limits those claims while leaving
the public/generated proposal work executable.

\section{The next complete qualification}

The remaining engineering route is concrete: choose source-audited packages
and permitted integration forms, implement representative bridges, and run
matched consumer workflows across all three surfaces. The trial must include
unsupported methods, late observations, graph ambiguity, weak or disconnected
effect sets, interruption and changed-input rejection. Passing results should
name the actual supported environments and domains.

This work will close an engineering gap only in those domains. Negative
persistent-map, latent-graph and dependent-learner studies remain scientific
obligations. Independent review, authentic bank validation and measured
economics remain additional gates. A coherent runtime makes their required
inputs and outputs reviewable; it does not turn missing information into an
accepted causal answer.

\label{exec13:operations-end}


% SOURCE docs/release/2026-10-29/source/docs/writing-repair/2026-10-29/src/91c37981ab684f02.tex CHAR 0
\section{Where a result becomes public after interruption}
\label{exec17:i4-publication}

A recovery specification should identify the transition at which a consumer
may rely on a result. The calculation finishing, the artifact bytes appearing,
and the job becoming completed are distinct events. In the inspected local
report worker, a job lock protects the final cancellation check, artifact
publication and update of the job's completed state and result reference.
The artifact store has its own content and metadata publications. This
source structure gives a concrete basis for reasoning about selected crash
points; it does not create a transaction over every file or an external
service \citep{exec13_ops_source}.

Consider the sequence $C\to A\to T$, where $C$ is a validated completed-work
checkpoint, $A$ means the report artifact has been stored, and $T$ means the
job record exposes a completed state and the artifact reference. Before $C$,
uncommitted work may need recalculation. Between $C$ and $A$, recovery can
reconstruct a report from the checkpoint after revalidating identities.
Between $A$ and $T$, an artifact may exist without being exposed by a
completed job. After $T$, a client that loses the response should retrieve
the result associated with that identity rather than create a new analysis.
These states explain why counting worker starts does not count distinct
scientific results, and why the existence of a report file alone should not
be the consumer's completion criterion.

The actual report path uses deterministic report contents and a content-based
artifact identity. A repeated publication checks the same content rather
than appending a second result to the public report. That supports recovery
for this path. A future estimator with stochastic or environment-sensitive
output would need a preserved randomization and numerical identity before
the same argument could apply. A new draw after a crash may be a new
statistical replicate; it is not automatically a continuation of the original
one. Likewise, byte-identical output can still embody false causal premises.
Recovery identity and scientific validity answer different questions.

Cancellation can be analyzed by the order of the protected transitions. In a
small authored state machine, begin at running. If cancellation is accepted
first, the state becomes cancelled and publication is disallowed. If the
publication transition completes first, the state becomes completed with a
public result and a later cancellation leaves that state intact. Enumeration
of these two orders reproduces the two outcomes
\citep{exec17_integration_review_reference}. The example assumes an indivisible
protected transition and a shared authoritative state. It is a reference
model for the local lock boundary, not a proof for independent machines,
uncooperative numerical kernels or external banking actions.

The lock also clarifies what a consumer can reasonably be told. A cancellation
request can wait while publication holds the lock. The returned completed
state then reports what actually happened; it must not claim that an earlier
user gesture erased the report. Conversely, a process killed while holding
the lock can leave an artifact published before the terminal state is
written. Recovery must recheck the job and artifact identities rather than
infer completion from the interrupted caller's intention. The preserved
fault fixtures exercise selected local interruptions, but they do not
enumerate every instruction boundary or prove crash consistency for all
filesystems and hardware \citep{exec13_ops_faults}.

A complete prototype fault protocol should state its injection points and
the invariants inspected after each restart. For the report path those
include: the same request and source identity; a valid completed-work prefix;
one contribution for each committed work identity; a valid artifact reference
only in a completed public state; and a truthful terminal cancellation
outcome. A negative case should change an input, checkpoint or environment
identity and require a refusal or new analysis rather than continuation.
These are proposal acceptance conditions. A platform study must additionally
state what durability its storage actually guarantees when the whole host,
not merely one worker, fails.

\subsection{Idempotence does not cross an uncooperative boundary}

An external offer illustrates the limit. Suppose an adapter sends a customer
message and loses the acknowledgement before recording success. There are
two worlds consistent with the local observation: the message was delivered,
or it was not. Blind retry duplicates the first world's action; refusing to
retry omits the second world's action. A local content hash cannot distinguish
them. A cooperating recipient that recognizes a durable action identity,
or an authoritative query that resolves delivery, can change the information
available. Without such a contract, the correct state is unresolved and a
reconciliation owner is needed. The current report service performs no such
action, so this example motivates a later interface requirement rather than
alleging an observed product failure.

This distinction should affect product scope and cost. A read-oriented report
prototype can demonstrate semantic preservation and local recovery without
having permission to act on customer accounts. An action-taking extension
adds authorization, recipient deduplication, reconciliation and compensation
requirements that must be costed and tested independently. A compensating
action, when one exists, is a new action with its own consequences; it is
not evidence that the original event never occurred. The proposal should
retain these boundaries when comparing workflow automation alternatives.

\section{Retention is reachability under a stable ownership decision}
\label{exec17:i4-retention}

Immutable artifacts make recovery auditable but create a storage obligation.
Deleting every file older than a chosen age can remove evidence still needed
by a current report or a resumed job. Keeping every file forever avoids that
particular mistake but makes cost and capacity unbounded in an ongoing
service. A proposed retention design should therefore distinguish the
scientific dependency graph from the policy that designates live roots.
The former records which exact objects are needed; the latter states which
reports, jobs, reviews and retention holds must remain available.

Let $\mathcal A$ be the finite set of stored artifact identities and let
$a\to b$ mean that artifact $a$ depends on $b$. Given a root set
$\mathcal R$, define the retained closure recursively by
\begin{equation}
 K_0=\mathcal R,\qquad
 K_{j+1}=K_j\cup\{b:\exists a\in K_j,\ a\to b\},\qquad
 K=\bigcup_{j\geq0}K_j.
 \label{exec17:i4-reachability}
\end{equation}
In a finite graph this process stabilizes after at most $|\mathcal A|$
strict additions of new identities, and cycles do not prevent termination
when visited identities are recorded. The deletion candidates are
$\mathcal A\setminus K$. This is a graph statement under complete dependency
metadata. It does not decide legal retention, scientific adequacy, access
authorization, or whether an owner has supplied every required root.

The exact reference has live report and checkpoint roots. Both reach a result
that depends on data and code. The retained set therefore contains those
five identities. An orphan artifact and an old result with its old input
form three deletion candidates when no other root references them.
The reference traverses the graph and checks this result, without deleting
any files \citep{exec17_integration_review_reference}. Adding an archived
review as a root could legitimately retain the old result. Removing an
unfinished job from the roots merely because its process is absent could
destroy recovery evidence. Root ownership is consequently a substantive
operating decision, not an incidental garbage-collector parameter.

Concurrent publication complicates the simple closure. A collector can
observe an artifact before a new job record points to it, classify it as
unreachable, and delete it just as the publisher exposes the reference.
A safe proposed implementation needs a stable root snapshot and coordination
with publication, or another protocol that prevents such a race. A
quarantine interval alone has no proof without an upper bound on publication
and reconciliation delays. The prototype can compare a lock-based local
collection phase with a versioned snapshot and explicit publication pins.
Each alternative should state the failure model and verify that a reachable
artifact is never deleted under that model.

The dependency graph should also preserve descriptor identities and custody.
A report referencing a small summary may still depend on a raw input that
must remain in a restricted store. Reachability does not require making all
reachable content readable to the same client. The proposed collector
operates under the store owner's authority, while a reader still passes
role, size and custody checks. Similarly, duplicate content can be stored
once physically while retaining distinct provenance or policy descriptors;
equating physical deduplication with interchangeable scientific evidence
would discard information about use and interpretation.

For a transparent storage calculation, let $n$ be admitted requests per day,
$s$ the mean additional retained bytes per request under a specified policy,
and $h$ the number of retained days. Ignoring shared content, metadata,
temporary files and replication gives $nsh$ bytes. This is an accounting
baseline, not a capacity guarantee. A real measurement must add those
categories and account for distributions: a few large covariance artifacts
can dominate a mean-sized planning example. Retention holds and historical
reviews can also extend effective lifetime beyond $h$. The proposal should
show which quantities are measured locally, stipulated for planning, or
awaiting a target operator's requirements.

\section{A worker limit is not an admission or deadline guarantee}
\label{exec17:i4-capacity}

The inspected worker has a memory and wall-time limit, while its coordinator
and SDK consume additional resources. An operator therefore needs a
whole-workflow admission rule. Suppose a hypothetical host budget is
1024 MiB, the resident parent requires 128 MiB, a reserve takes another
128 MiB, and each admitted worker is allotted 256 MiB. Under an additive
worst-case model the maximum concurrent count is
\begin{equation}
 k=\left\lfloor\frac{1024-128-128}{256}\right\rfloor=3.
 \label{exec17:i4-memory}
\end{equation}
Four workers would consume the entire 1024 MiB host budget, leaving no allowance for either the parent or the reserve. These
figures are deliberately stipulated and reproduced by the reference
\citep{exec17_integration_review_reference}. Actual address-space limits are
not measurements of peak resident memory, and shared pages, buffers,
serialization and numerical libraries can alter the relationship. The
equation describes a conservative admission model once its quantities and
scope have been justified; it does not qualify the present host.

Admission also affects delay before work starts. Consider five simultaneous
requests, two identical workers, and exactly three seconds of service per
request. With a first-available-worker rule, completion times are
$3,3,6,6,9$ seconds. A deadline eight seconds after arrival is met by four
requests and missed by one, even though every service lasts only three
seconds. Direct scheduling and the independent expression
$3\lceil j/2\rceil$, $j=1,\ldots,5$, agree in the exact reference.
Thus a per-worker time limit can hold while the user-facing deadline fails.
The denominator includes the fifth request. Dropping it as an unfinished
job would overstate deadline performance.

The relevant deadline must be declared at intake. If $t_i$ is arrival time
and $b_i$ is the allowed end-to-end duration, a later start at $s_i$ leaves
at most $t_i+b_i-s_i$ for remaining work, less any required publication and
delivery allowance. Starting a fresh full budget after queueing changes the
contract. If the remaining allowance is already exhausted, rejecting or
expiring the request is more truthful than launching work and relabeling
its eventual result on-time. If the service-time bound is unknown, a queue
policy can report an estimate with its limitations; it cannot provide a
hard completion promise merely by limiting concurrency.

The queue example intentionally removes random arrivals, heterogeneous
service times and failures. Those extensions need a stated workload law or
measured trace, not a casual independence assumption. Correlated daily
campaigns can create bursts even when the long-run average arrival rate
looks modest. Long jobs can delay short ones under one policy; prioritizing
short jobs can systematically exclude the scientifically difficult tasks.
The proposal should compare admission and scheduling rules using complete
task identities and the same mix of targets, with rejected, expired and
cancelled requests visible. A faster successful subset is not a generally
faster workflow.

Back-pressure is therefore a scientific-interface issue as well as an
operating concern. A rejected request needs a distinct execution state and
its original target. It must not become an unsupported causal-method claim.
A timeout during a difficult latent-graph calculation cannot be counted as
proof of nonidentification. Conversely, an identified but unbounded result
should not be retried until it becomes a convenient scalar. Capacity,
scientific disposition and usefulness should be recorded in separate fields
so that each failure has an appropriate remedy.

\subsection{What an inexpensive complete trial should measure}

A public/generated prototype trial can now be specified without pretending
to have production data. Freeze a finite mixture of supported and negative
scientific requests, artifact sizes and budget classes. Record intake,
admission, queue entry, worker start, computation completion, publication
and client receipt under one clock convention. Measure peak resources for
the process tree and store, include failed parsing and serialization, and
inject declared crashes before and after checkpoint and publication. The
retention trial should operate on a disposable copy with known dependencies
and verify every live reference after collection. The reference constructions
in this section make expected states auditable; they do not substitute for
the subsequent measurements.

Compare the library, CLI and compatible MCP path under the same request
definitions, while separately reporting their unavoidable launch and
transport costs. Use an incumbent analyst or scripted workflow as a
decision comparator only when its task and permitted information match.
Monetary costs should distinguish one-time construction, shared operation,
per-request compute, review, retry and recovery. Avoid counting a worker's
CPU charge again inside an already inclusive service charge. The prior
all-attempt example shows why costs from failures and fallback decisions
belong in the economic denominator \citep{exec13_ops_exact}.

The resulting evidence can support a scoped design decision: which consumer
boundary preserves meaning, which storage rule maintains recovery, and what
workload the local prototype can actually process. It cannot authenticate a
bank owner, supply real assignment truth or establish institution-scale
availability. Those later trials remain conditional on prototype merit and
agreed acceptance criteria. Maintaining that distinction permits the present
proposal to finish its engineering argument while keeping the expensive
external commitments explicit.
\label{exec17:i4-capacity-end}


% SOURCE docs/release/2026-10-29/source/docs/writing-repair/2026-10-29/src/4109f6038aec3b0a.tex CHAR 0
\section{A retention protocol that survives a publishing race}
\label{exec17:i4-pins}
The preceding reachability rule states which artifacts are needed at one instant. A collector must additionally preserve that property while a report becomes visible. Consider a generated deposit report whose immutable dependency descriptor points to the input balance extract. Before publication, neither object need be reachable from an accepted report. A publisher can be preparing a legitimate result precisely when a collector takes its snapshot. The substantive operating question is whether the proposed coordination rule prevents a published report from losing its input, without assuming that every process which started will eventually finish.

For a finite model, let $V$ be the set of present immutable object identities, $R$ the accepted report roots, and $P$ the roots pinned by admitted but unfinished publishers. Each descriptor has an immutable, complete list of dependencies. Write $\operatorname{cl}(S)$ for $S$ together with all objects reachable through these lists. The operating invariant is
\begin{equation}
 \operatorname{cl}(R\cup P)\subseteq V.
 \label{exec17:i4-live-invariant}
\end{equation}
This is a storage-availability condition for declared dependencies. It cannot verify that a descriptor lists all the scientific inputs actually used, that a balance extract is accurate, or that a customer record may be retained. Those questions require separate provenance and custody evidence. The following protocol is a proposed local design and finite reference, not an implemented collector or a theorem about arbitrary distributed storage~\citep{exec17_retention_reference}.

\subsection{Pin acquisition must agree with the deletion decision}
An admitted publisher needs a reservation that protects a complete dependency closure. Under the proposed local rule, pin acquisition checks that all required objects exist and adds the pin while holding the same metadata coordination boundary used by deletion. The publisher subsequently installs the report root before releasing its pin. Publication, pin release and the collector's final deletion decision use that boundary too. An object being staged outside the store is not yet promised to any consumer. If collection wins before admission, the publisher must reconstruct or reacquire the missing objects and seek admission again; it cannot claim a valid pin on already missing bytes.

Let $M$ be the closure marked at the collector's earlier snapshot. Immediately before deletion, re-read the current roots and pins under coordination and remove only
\begin{equation}
 D=V\setminus\bigl[M\cup\operatorname{cl}(R\cup P)\bigr].
 \label{exec17:i4-delete-set}
\end{equation}
An old mark can keep extra objects, which costs storage, but the current closure prevents a newly protected object from being collected. The final dependency check and deletion must be indivisible with respect to pin acquisition and root publication in this model. Checking first, releasing the lock, and deleting later reintroduces the same race. A scalable implementation may use another mechanism, but it owes an equivalent argument; naming a snapshot or a quarantine does not establish that argument.

The invariant has a short inductive proof. Begin with all current roots and pins present. A successful pin acquisition checks its entire closure, so adding that pin preserves the inclusion. Rejected admission changes neither roots nor pins. Publication adds a root whose closure was protected by a pin, and occurs before that pin is removed. Releasing a pin can only reduce the left side of the inclusion. A snapshot changes no live object. Finally, deletion by the displayed rule removes nothing in the current closure. Every permitted transition therefore preserves the invariant. This proof depends on complete immutable descriptors and the shared coordination boundary; it does not follow merely from individually atomic file renames.

\subsection{Ten schedules separate safe rejection from broken publication}
The reference has two objects: a report and its input, with one edge from the report to the input. A publisher performs pin, publish and unpin in that order. A collector performs snapshot then sweep. There are $\binom{5}{3}=10$ order-preserving interleavings. This denominator follows by choosing the three publisher positions; it does not assume equally likely runtime schedules. The reference enumerates all ten and retains the state after each event~\citep{exec17_retention_reference}.

Three interleavings violate the invariant when sweep uses only its stale snapshot. One witness is snapshot, pin, publish, sweep, unpin. The snapshot sees no roots. Pin acquisition then sees both objects, and publication exposes the report. A stale sweep deletes both objects despite that successful publication. A second order, snapshot, pin, sweep, publish, unpin, breaks the pin before publication. In the third, snapshot, pin, publish, unpin, sweep, the published root still fails to protect its objects because the collector never rechecks it. The defect is independent of whether the numerical report was computed correctly.

With the current-closure recheck, all ten schedules preserve the invariant. Nine publish successfully. In the remaining order, snapshot and sweep both occur before pin acquisition; both unprotected objects have been collected, so admission rejects and no report becomes public. This is an availability failure for that attempt, yet it is the correct integrity outcome. Treating a rejected attempt as an unsafe publication would confound two distinct criteria. Conversely, counting only the nine successful attempts and suppressing the rejection would overstate availability. A subsequent implementation trial should retain both denominators.

The example does not estimate a crash probability or report a measured three-in-ten failure rate. It is exhaustive only for this five-event model. Real implementations add content writes, descriptor writes, error returns, multiple publishers and filesystem failures. The proof identifies which transitions require a shared boundary and which additional events need refinement. The next low-cost experiment can therefore test a small disposable store with deliberate pauses at these boundaries, rather than run an expensive bank-data study to rediscover a deterministic race.

\subsection{A crashed publisher creates a liveness obligation}
Pins introduce a second question: who releases a pin after its publisher disappears? Leaving a stale pin in place protects its complete closure indefinitely. The reference's two-object example retains both objects, even if no accepted report will ever need them. This is safe under the inclusion invariant but can exhaust storage. Safety and eventual reclamation are separate requirements; a system should not claim both from the same reachability calculation.

Deleting a pin after a timeout is insufficient if a paused publisher can later resume. Once the expired pin is removed, a fresh sweep can collect its report and input. An old publisher which resumes and exposes the report without reacquiring admission then creates a dangling reference. The finite reference records exactly this counterexample. No numerical choice of timeout rules it out without an assumption bounding pauses or a mechanism preventing a stale publisher from committing~\citep{exec17_retention_reference}.

A proposed repair associates the publication right with a generation token. Revocation changes the current generation under coordination; publication must check that its token remains current and that its closure remains present. An old publisher can finish computing, but cannot expose its result using a revoked token. It may start a new admission attempt and reuse still-valid artifacts. This proposal needs a precise failure model before implementation: the generation check and root publication must not themselves be separable by revocation. It is not sufficient to check the token once when a long computation starts.

A local prototype can begin with a simpler stop-the-world collection phase: prevent new admissions, reconcile existing jobs, collect under one owner, then reopen admission. That sacrifices availability but makes the boundary easier to inspect. A concurrent design with publication pins permits more overlap and requires stronger coordination and recovery logic. The choice should depend on the expected workflow and tolerable interruption, not a general preference for concurrency. Neither alternative licenses deletion of records subject to an external retention hold.

\subsection{Retention classes must remain separate from readable resources}
An accepted report root, an unfinished job pin, a reproducibility hold and a custodial hold have different owners and release conditions. The storage design should record which class is retaining each object. A completed report may no longer be exposed to an agent while a review hold still requires its inputs. Conversely, a visible summary does not authorize the agent to read every reachable raw record. Closure determines what must remain available to authorized recovery; it does not grant a new reader permission.

This distinction matters in a deposit analysis that joins a restricted customer extract to an openly distributed method description. Physical deduplication of identical bytes may save space while their custody descriptors remain different. A release condition on one descriptor cannot silently release a hold belonging to another owner. The collector should compute the union of admissible retention obligations, then apply store-specific deletion authority. The proposal cannot specify actual institutional retention periods or assert that indefinite retention is lawful. Those are future owner inputs. Its current contribution is to expose the interfaces at which such inputs change the deletion decision.

Restoring a backup presents the dual problem. Recovered report bytes are insufficient if their descriptors point to absent input or method identities. Before making restored roots visible, compute the dependency closure against the restored inventory, validate the recorded identities and classify missing dependencies. A report with a missing input can be retained as an incomplete historical artifact while remaining ineligible for a reproducibility claim. Recomputing from a newer extract creates a successor analysis, not restoration of the old result. A recovery test should deliberately omit one dependency and verify this distinction.

The resulting prototype protocol is concrete. Freeze an immutable dependency fixture, declare the root and pin lifecycle, enumerate or inject the publication/collection interleavings, interrupt a publisher before pin release, attempt stale-token publication, and restore a partial inventory. Record every admission, rejection, published root and missing dependency. Measure retained bytes and collection delay separately from integrity violations. The finite reference supplies exact expectations for its small core; the added token, custody and backup behaviors still require implementation and their own checks. Such public/generated evidence can establish whether the proposed operating boundary is useful before seeking authentic platform owners and costly banking experiments.
\label{exec17:i4-pins-end}


% SOURCE docs/release/2026-10-29/source/docs/writing-repair/2026-10-29/src/633c351e4ada200e.tex CHAR 0
\section{A discovery journey that preserves refusal and uncertainty}
\label{exec13:discovery-journey}

A bank's request to discover why an offer changes twelve-month customer value
contains several operations. The information must precede the offer, the
candidate graph must retain its unresolved orientations, the requested effect
must be identified under a declared model, and the financial interpretation
must preserve the resulting uncertainty. A common command-line or library
interface is useful only if its callers can tell which of these operations
actually occurred.

The local reference service supports graph syntax review and a particular
supplied-DAG back-door calculation. Its \texttt{graph.discover} request was unsupported in the recorded exercise. This observation binds that service path and revision, rather than every discovery construction developed elsewhere in the book. A syntax response can preserve a CPDAG or PAG and its endpoint
marks without establishing that the graph is true or recognizing every valid
member of the represented class. The supplied-DAG identification method does
not thereby become a method for arbitrary graph classes, latent confounding,
temporal discovery or dependent sampling.

The extended local exercise froze 24 authored cases and 98 invocation slots
before executing them~\citep{exec13_discovery_journey}. Each case was called
through the library and CLI, initially and after restart or re-instantiation.
Two additional requests retained a completed request's key while changing its
target horizon. The elementary references were graph paths, factored
inequalities, exact timestamp order and explicit changes to request identity.
They support a deterministic behavior comparison, not a statistical discovery
accuracy or causal-coverage estimate.

\subsection{Identification and refusal in a concrete offer graph}

For the supplied graph $A\to Y$, an empty adjustment set satisfies the selected
back-door criterion. Add a pre-offer risk variable $R$ with
$A\leftarrow R\to Y$. If $R$ is observed and included in the adjustment
request, the method returns the corresponding conditional identification
expression. Omitting $R$ leaves the path open. Declaring it latent does not
make an observed adjustment possible. Conversely, requesting adjustment for
the mediator in $A\to R\to Y$ fails this particular criterion for the total
effect. Such failures must remain failures of the attempted method; they are
not proofs that every identification strategy is impossible.

The cases also require refusal when a declared premise is contradicted,
unsupported dispositions for CPDAG/PAG identification and dependent-month
sampling, and an invalid-request response for a directed cycle. Correctly
recognizing these boundaries is useful behavior. It does not supply the causal
effect that the bank originally wanted.

\subsection{Numerical sets, information clocks and immutable context}

Four deterministic sensitivity calculations check set serialization:
$1-x^2\le0$ gives two rays, $1+x^2\le0$ gives the empty set,
$-1-x^2\le0$ gives the whole line, and $(x-1)^2\le0$ gives a singleton.
The service retains those different topologies. The examples have no sampling
coverage probability; numerical edge cases outside these selected polynomials
still need an independent error analysis.

The observation audit admits the timestamp row satisfying
$\text{event}\le\text{available}<\text{decision}$. Availability exactly at
the decision, a missing timestamp, and an availability time preceding the event
are retained as failures. These declarations alone cannot establish that a
real source system recorded the correct timestamps.

Two artifact cases test a different failure mechanism. Changing the sign of
the target after importing a moment artifact must prevent its reuse under the
new scientific context. Corrupting the stored bytes must fail their identity
check. Keeping an idempotency key while changing the horizon must yield a
conflict. A client that receives such a failure has no new computed effect to
display. Its request correlation ID can change across a replay, while the
scientific content must remain equal.

\subsection{Observed behavior and the cost of this local exercise}

All 98 slots matched their frozen expectations, and all 24 case comparisons
retained the same scientific responses across the four route/restart calls.
There were 44 computed responses, 36 unsupported responses, four refusals,
four invalid requests and ten unavailable responses, including the two target
conflicts. The computed group includes capabilities, syntax, conditional
expressions, sensitivity sets and an audit; it is not a count of 44 estimated
causal effects. All groups remain in the denominator.

The run resumed after 40 and 80 recorded invocations, preserving both completed
prefixes exactly. The 98 invocations took approximately 41.99 seconds in this
local run, including worker/process launch, failed requests and replay. This
measurement excludes author effort, process startup before the runner's timer,
and pauses between runner segments. Separate active-runner time includes
artifact imports and journaling. Reported resident-memory readings are maxima
for individual parent/child processes, not concurrent process-tree memory.
These measurements neither establish bank throughput nor estimate bank ROI.

Restarting these calls is a narrower exercise than killing a multi-host
service mid-operation. A crash between a completed call and its journal write
can repeat that call with the same key; the exercise does not establish
exactly-once execution. It also runs no actual sibling consumer, MCP transport,
live language model or installed discovery engine. Those distinctions preserve
the meaning of both positive checks and remaining work.

\subsection{How the proposed consumer uses the result}

A consumer first reads capabilities and checks its clock, target, sampling law,
graph type and requested method. A syntax pass allows a graph to be displayed
with its marks intact. A conditional identification expression permits the
next estimation design to be considered. An unsupported or refused result
routes the question to another justified method or requests missing information.
It must not trigger an automatic point estimate from a language model.

A contract double can demonstrate these client decisions using typed local
responses. Acceptance of an actual workspace integration additionally needs
the owner's request and unit map, source clock, execution trace, failure and
restart results, and an explicit disposition. This is later evidence about
that integration. For the present proposal, the local exercise already makes
the required behavior concrete and reproducible while preserving the missing
discovery-to-effect capabilities.
\label{exec13:discovery-journey-end}


% SOURCE docs/release/2026-10-29/source/docs/writing-repair/2026-10-29/src/6e8306b8cc42e7f2.tex CHAR 0
\section{From a funding response to a defensible treasury decision}
\label{exec13:funding-journey}

A treasury team considering a deposit-rate change needs more than a response
coefficient. It needs to know which balance the response concerns, whose
accounts constitute the target population, which information was available
when the quote was chosen, and whether the uncertainty survives the journey
from estimated scores to a financial recommendation. This section develops
that journey through exact project examples and a bounded local interface
exercise. The examples are generated constructions dated 13 October 2026;
they are not observations of a bank or later additions to the literature
search cutoff~\citep{exec13_funding_exact,exec13_funding_runtime}.

Let $\beta$ denote the response of the next twelve months' \emph{average}
retained deposit balance, measured in USD millions, to a one-percentage-point
increase in the quoted deposit rate. A response of the balance at month twelve
is a different quantity: earning a spread on a closing balance does not
describe the interest income earned over the preceding year. A
quarter-percentage-point action has response $0.25\beta$ only under an
additional local-linearity and intervention-stability premise. A normalized
proxy response does not by itself establish that a feasible deposit-rate
decision implements the corresponding structural intervention.

These distinctions belong in the proposed product's scientific context.
The target, horizon, units, sign, intervention and sampling law must travel
with the estimate. A consumer also needs evidence for the population map,
inclusion mechanism, score construction and covariance. It should be possible
to compute a conditional mathematical object while withholding a funding
recommendation whose premises remain unresolved. The existing service
implements the former operation on supplied summaries; the broader decision
layer remains proposed~\citep{exec13_funding_service}.

\subsection{Information clocks and the population behind the balance}

A feature available to a quote decision must satisfy
\[
 t_{\rm event}\le t_{\rm available}<t_{\rm decision}.
\]
The first time that a value could actually be used is the relevant
availability time. Substituting a subsequently revised figure into a past
feature record can violate this requirement even if the revision carries
the same economic reference date. The rule concerns decision information.
Future outcomes are necessarily observed later and instead require an
explicit follow-up and censoring rule. An audit that rejected every
post-decision outcome would be auditing the wrong object.

Correct timestamp order cannot establish that two records belong to the same
household. To see why this matters, consider four generated households with
truth-known contributions $(0,2,4,6)$. Their household mean is $3$. Repeating
the first household's zero contribution in four records gives seven records,
\[
 (0,0,0,0,2,4,6),\qquad \overline{y}_{\rm record}=12/7.
\]
The true household mean remains $3$ when the four zero records have a common
identity. If they represent four distinct households, $12/7$ is instead the
mean of that seven-household population. Identical values and valid clocks
cannot distinguish these two interpretations. An artifact hash establishes
which bytes were used; it does not recover the missing unit map.

Inclusion creates a separate problem even with known identities. Suppose the
four households are independently included with probabilities
$(1,1,1/2,1/2)$. The first two are always present; the last two generate four
equally probable selection states. The observed unweighted means are
\[
 1,\quad 2,\quad 8/3,\quad 3,
 \qquad
 E[\overline{y}_{\rm observed}]=13/6.
\]
For known target size $N=4$ and known positive probabilities $\pi_i$, the
population-total estimator divided by $N$ is
\[
 \widehat\mu_{\rm HT}=\frac{1}{4}\sum_{i=1}^{4}
       \frac{S_i y_i}{\pi_i},
 \qquad
 E[\widehat\mu_{\rm HT}]=3.
\]
Its four possible values are $1/2,5/2,7/2,11/2$. The equality follows by
enumeration or by $E[S_i/\pi_i]=1$; it does not follow from observing only
the selected records. Unknown inclusion probabilities, an unknown target
population or a stratum with zero inclusion require additional information
or a restricted target. The contributions in this construction are generated
truth. The calculation does not make both potential outcomes observable in
a real funding study~\citep{exec13_funding_exact}.

For the proposed workflow, unit and inclusion evidence therefore precede
claims about a target-bank response. A useful public-panel prototype can
exercise institution identifiers, dates and observed balance aggregation.
It should not label public institution identifiers as a verified household
map, nor claim assignment truth from a panel that does not record it.
Generated studies remain useful because the hidden truth can be deliberately
corrupted and kept out of the estimation inputs.

\subsection{The covariance must include what the workflow estimated}

Write a normalized response as $\beta=\mu_Y/\mu_D$, where $\mu_Y$ is a
numerator moment and $\mu_D$ is the normalization moment. The estimator may
first fit a conditional mean, propensity, proxy, linkage model or other
nuisance quantity. If a justified expansion has leading score
\[
 \psi_t=\psi_{0,t}+A\phi_{\eta,t},
\]
then, when the corresponding stationary long-run covariance exists, it is
\begin{align*}
 \Omega
 &=\sum_{k=-\infty}^{\infty}
   \operatorname{Cov}(\psi_t,\psi_{t-k})\\
 &=\sum_k\bigl\{
   \operatorname{Cov}(\psi_{0,t},\psi_{0,t-k})
   +\operatorname{Cov}(\psi_{0,t},\phi_{\eta,t-k})A^{\mathsf T}
   \\
 &\qquad+A\operatorname{Cov}(\phi_{\eta,t},\psi_{0,t-k})
   +A\operatorname{Cov}(\phi_{\eta,t},\phi_{\eta,t-k})A^{\mathsf T}
   \bigr\}.
\end{align*}
Here $A$ maps the nuisance influence into the moment coordinates. This is an
algebraic consequence of a supplied expansion, not a proof that an arbitrary
learner has one. Overlapping twelve-month outcomes can introduce dependence
across monthly origins. Cross-outcome and denominator covariance must also
survive aggregation. A label saying ``dependent monthly origins'' does not
establish that the supplied matrix includes these contributions.

An exact Gaussian example makes the possible omission concrete. Let
\[
 X_i=\theta-2\eta_0+U_i,\qquad
 \widehat\eta=\eta_0+\overline V,\qquad
 \widehat\theta=\overline X+2\widehat\eta,
\]
where $U_i\sim N(0,1)$ and $V_i\sim N(0,2)$ are independent within and
across observations. Then
\[
 \widehat\theta-\theta=\overline U+2\overline V,
 \qquad n\operatorname{Var}(\widehat\theta)=1+4(2)=9.
\]
Retaining only the $U$ contribution supplies variance $1$. At $n=400$ and
normal critical value $z=2$, an illustrative estimate of $2$ receives intervals
$[1.9,2.1]$ and $[1.7,2.3]$, respectively. In this stipulated Gaussian law,
the variance-nine interval has exact coverage $2\Phi(2)-1\approx0.95450$.
The variance-one interval's coverage is instead
$2\Phi(2/3)-1\approx0.49501$. These are analytic probabilities under a known
law, not estimates of empirical performance. The example is intentionally
nonorthogonal and independent; it does not settle the dependent or
arbitrary-learner problem~\citep{exec13_funding_exact}.

Orthogonal scores, suitable sample splitting or an explicit joint
influence-function construction are substantive alternatives. Each needs
its own assumptions and remainder control. With dependence, a nominal
split of adjacent months may retain information leakage and correlation.
A proposal should specify which learner, split rule, dependence class and
variance construction a minimal prototype will actually support. It should
then compare that restricted path with a matched simpler baseline and retain
the regimes in which the argument or precision fails.

\subsection{Invert the residual; preserve the resulting set}

For supplied moment estimates $(\widehat\mu_Y,\widehat\mu_D)$, sample
size $n$, covariance $\widehat\Omega$ and normal critical value $z$,
the local calculation retains candidate $b$ satisfying
\begin{equation}
 n(\widehat\mu_Y-b\widehat\mu_D)^2
 \le z^2\bigl(
   \widehat\Omega_{YY}-2b\widehat\Omega_{YD}
      +b^2\widehat\Omega_{DD}\bigr).
 \label{exec13:funding-inversion}
\end{equation}
The inequality avoids numerical division by a small normalization estimate.
Expanding its two sides yields a quadratic sublevel set. Statistical coverage
still requires a justified moment approximation and covariance. In
particular, a serialized positive-semidefinite matrix is not proof of its
sampling validity~\citep{exec13_funding_service}.

The following examples use $z=2$. The one-coordinate error probability is
$\alpha=2\{1-\Phi(2)\}\approx0.0455003$.
\begin{center}
\small
\begin{tabular}{p{.31\linewidth}p{.28\linewidth}p{.29\linewidth}}
\toprule
Supplied moments and covariance & Residual inequality & Retained set\\
\midrule
$n=100$, means $(2,1)$, covariance $I_2$
& $96b^2-400b+396\le0$
& $[(25-\sqrt{31})/12,(25+\sqrt{31})/12]$\\[5pt]
$n=4$, means $(1,0)$, covariance $\operatorname{diag}(0,1)$
& $4-4b^2\le0$
& $(-\infty,-1]\cup[1,\infty)$\\[5pt]
$n=4$, means $(0,0)$, covariance $I_2$
& $-4(1+b^2)\le0$
& $\mathbb R$\\[5pt]
$n=4$, means $(1,0)$, covariance zero
& $4\le0$
& $\varnothing$\\[5pt]
$n=4$, means $(2,1)$, covariance zero
& $4(b-2)^2\le0$
& $\{2\}$\\
\bottomrule
\end{tabular}
\end{center}
The zero-variance rows deliberately test mathematical and interface
boundaries. They are not valid evidence of perfect precision in a financial
application. An empty set can indicate incompatible inputs or a failed
model; it must not be interpreted as proof that all financial risks vanish.
Two rays cannot be replaced by their midpoint or by a finite interval without
changing the result. The whole real line records a lack of useful restriction.

For several outcomes sharing a denominator, the service can use a
Bonferroni critical value for each coordinate. The two-outcome example has
means $(2,1,1)$ and supplied covariance
\[
 \begin{pmatrix}1&1/2&0\\1/2&4&0\\0&0&0\end{pmatrix}.
\]
With $n=400$ and family error probability $2\alpha$, the coordinate intervals
are $[1.9,2.1]$ and $[0.8,1.2]$. Their labels preserve the common construction
and its nominal family level $1-2\alpha\approx0.90900$. This is not a
95\% simultaneous result. A union-bound argument, when individual coverage
is justified, can apply despite cross-coordinate correlation; it does not
make the resulting rectangle an efficient fitted joint region. The transport
field named \texttt{joint\_confidence} must be read with that construction
and its supplied-covariance qualification.

\subsection{Useful precision is a property of the decision}

Consider an explicitly assumed annual funding scenario. A baseline average
balance of USD 2.6 million is repriced by 0.25 percentage points, costing
USD 6,500. Incremental retained funds earn a 1.5\% annual spread
\emph{net of their own offered deposit interest}. The incremental average
balance is $0.25\beta$ million USD. With an illustrative workflow allowance
of USD 500, net profit is
\begin{equation}
 \Pi(\beta)
 =10^6(0.25\beta)(0.015)
   -2.6\mathbin{\cdot}10^6(0.0025)-500
 =3750\beta-7000\quad\text{USD}.
 \label{exec13:funding-profit}
\end{equation}
The average-balance convention is necessary for this annual spread
calculation. A different repricing base, within-year balance path, liquidity
cost or net spread changes the function and should be represented explicitly.

The scenario allowance assigns USD 120 to data and identity preparation,
150 to integration amortization, 100 to scientific review, 30 to computation
and 100 to monitoring and recovery. These are planning assumptions. Actual
human adjudication, optional model calls, platform review, implementation,
security, deployment and continuing operations may exceed them. Failed
requests and retries belong in the workflow denominator. The allowance
cannot be presented as an observed bank price.

The break-even response is $\beta=28/15$. A point estimate $\beta=2$
gives profit USD 500. Mapping the two intervals from the fitted-nuisance
example through the increasing function in
\eqref{exec13:funding-profit} gives
\[
 \Pi([1.9,2.1])=[125,875],\qquad
 \Pi([1.7,2.3])=[-625,1625]\quad\text{USD}.
\]
The first interval excludes a loss under the stated cost scenario; the
second does not. This is the practical importance of the missing variance
component. It neither demonstrates bank profit nor makes every interval
decision-ready. The identification, unit map and covariance premises must
also be justified. Uncertainty about costs can be propagated by taking the
financial image over their declared scenario set as well.

For the weak-denominator rays, a finite favorable point estimate is especially
misleading. Under the same monotone profit map their image is
\[
 (-\infty,-10750]\ \cup\ [-3250,\infty).
\]
Neither a finite lower bound nor a narrow range is available. Increasing the
number of observations need not resolve a structurally weak or invalid
normalization. Possible responses include changing the target, seeking a
better justified design or reporting sensitivity. Choosing a favorable
representative point does not resolve the information problem.

\subsection{What the local library and CLI actually demonstrated}

The bounded exercise specified 29 authored cases, each invoked through the
library and CLI initially and after re-instantiation or process restart.
Two further calls changed the horizon under a previously used key. The
resulting 118 service slots all matched their fixed expectations after two
disclosed harness repairs. The original harness first omitted a standard
library import and later referred to a nonexistent copied fixture identifier.
The first failure preceded any service call; the second occurred after 116
completed slots. Both failures and the unchanged 116-result prefix were
preserved, and the corrections were recorded before the affected calls.
This is evidence from an amended execution, not evidence that the original
executable ran without repair~\citep{exec13_funding_runtime}.

Of the 118 responses, 48 were computed outputs, 40 were failed or unavailable
input results, eight refused contradicted premises, twenty reported
unsupported operations and two rejected changed-horizon reuse. The computed
outputs include diagnostics, deterministic sensitivity, unbounded and empty
sets. They must not be counted as 48 useful causal estimates. All 29
library/CLI/restart comparisons agreed on complete semantic responses,
excluding only the request identifier. Independent decimal root calculations
and 451 direct rational-grid membership checks agreed with the retained
sets. This finite screen does not estimate statistical coverage or a
population failure probability.

The feature-clock case retained four failures: availability at the decision,
a missing availability field, availability before the event, and timestamps
without time zones. The valid row passed. Changing units or the availability
field after binding a moment artifact caused a context failure, and corrupting
the artifact bytes caused an access failure. Continuous-treatment and
dependent-unit back-door requests, identity/inclusion audits, nuisance fitting
and design evaluation remained unsupported in the selected paths.

Most consequentially, summaries with explicitly unresolved household identity
or unknown inclusion still returned conditional intervals. Contradicted
premises caused refusal, but an unverified premise was not automatically
treated as a contradiction. The service also retained a declared dependent
sampling label without verifying an upstream long-run covariance estimator.
These observed boundaries mean the proposed consumer needs an explicit
prerequisite review before presenting a financial recommendation. Passing
schema validation, a hash check or a clock audit is insufficient.

The measured sum of invocation times was approximately 56.46 seconds, including
failed, replayed and rejected service requests and their worker or CLI startup.
This is a local measurement of the stated exercise. The initial failed
harness startup was not timed, and the measure excludes authoring, repairs,
review, manuscript work and idle intervals. Process high-water memory is not
simultaneous memory for an entire process tree. These measurements therefore
cannot establish complete bank workflow cost or production throughput.
No new MCP transport qualification, foreign-language adapter, live model or
actual workspace consumer was exercised.

\subsection{A minimal next prototype with an informative stopping rule}

The next prototype should first restrict the scientific claim. One useful
design jointly generates dependent monthly funding histories, known
households, observed and missing follow-up, and a specified nuisance
estimator. It can compare covariance methods under the same generated law,
keeping the household truth and future outcomes outside the analysis inputs.
Useful outputs are all-attempt coverage where justified, set topology,
availability, width relative to the financial threshold, erroneous action
rates and complete resource categories. A matched simple estimator supplies
a baseline; a flexible learner must justify its extra cost and assumptions.

This experiment needs a reviewed derivation and prospectively fixed
regimes, replications, criteria and stopping rule. The finite conformance screen described here cannot supply those performance results. The fitted generated-panel study in Chapter~\ref{gap29:funding} and the generated linkage study in Chapter~\ref{gap29:observation} provide separate evidence for their own laws and targets. Neither converts this summary-input screen into qualification of the complete dependent-household design proposed here. Public funding
panels can inform plausible scales, release delays and institutional
continuity. Generated truth supplies controlled interventions and unit maps
that the public observations may lack. Neither source should be credited
with information that it does not contain.

A contract double for a treasury consumer can separately test whether absent
identity, inclusion, covariance or decision-cost evidence causes the intended
hold disposition. Such a double demonstrates a proposed consumer boundary;
it is not an accepted sibling-repository workflow. An owner-run bank study
becomes informative after a bounded prototype identifies a consequential
remaining uncertainty, an attainable precision target and plausible value
from resolving it. This preserves a practical path from proposal to evidence
while retaining the distinction between a computed response set and a
defensible funding decision.
\label{exec13:funding-journey-end}


% SOURCE docs/release/2026-10-29/source/docs/writing-repair/2026-10-29/src/146d5f4fb9fee477.tex CHAR 0
\section{An observation contract from source records to a contact decision}
\label{exec13:observation-journey}

A servicing manager wants to know whether contacting an eligible borrower
reduces the probability of an adverse event during the following twelve
months. Before comparing contact with no contact, the analysis must establish
which borrower each record describes, what the adverse-event label means,
which borrowers could enter the observed sample, and which information was
available when contact was chosen. These requirements are related, but one
cannot stand in for another. A correctly dated row can belong to the wrong
borrower; a correctly linked record can carry a noisy label; and a table of
correct labels can represent a selected population. The proposed observation
contract makes these distinctions explicit before downstream causal analysis.

The immediate design uses generated laws and the existing local library and
command-line service. Its purpose is to establish what a minimal prototype
can preserve and which further operation is justified. It does not require
an expensive bank validation merely to explain a failure mechanism. The
development sequence is to specify the observation law, demonstrate the
consequence of violating it, implement a narrow supported path, and then ask
whether an unresolved bank-specific question could change the investment
decision~\citep{exec13_observation_design}.

\subsection{Keep the information clock separate from outcome ascertainment}

For a feature record, write event time $t_e$, first usable availability
$t_a$, and decision time $t_d$. The selected admission rule is
$t_e\leq t_a<t_d$. Equality of event and availability is permitted; equality
of availability and decision is not. A corrected value published after the
decision retains its later availability even when it refers to an earlier
event. Comparing timezone-aware instants avoids inventing a time ordering
from differently displayed local times. None of these checks establishes
that the supplied timestamps themselves are truthful.

An outcome has a different role. An event six months after contact can be
an eligible twelve-month outcome, although it must not appear among the
features available before contact. Its ascertainment date, censoring rule,
revision history and horizon belong in the outcome contract. Applying the
feature clock to every row would reject legitimate follow-up observations;
ignoring availability would admit future information. The proposed consumer
therefore declares the row's role before selecting an audit. The present
reference service implements the strict feature clock only and explicitly
withholds its alternative-clock operation~\citep{exec13_observation_primary}.

The proposed manifest also retains the source identity and owner declaration,
license and version, extraction and transformation history, origin and
household identifiers, candidate linkage maps, inclusion declarations,
label definition and adjudication evidence, and the counts of accepted,
rejected and unavailable records. A byte hash binds the content that was
used. It cannot establish who originally observed an event or whether the
contents are correct. Source custody, content integrity and scientific
measurement validity consequently have separate evidence fields and separate
failure reasons~\citep{exec13_observation_design}.

\subsection{Marginal inclusion probabilities do not determine precision}

Consider a fixed population with two origin values $y_1=20$ and $y_2=80$.
The equally weighted population mean is 50. Let $R_i$ indicate inclusion
and suppose both marginal inclusion probabilities are $\pi_i=1/2$. The
expansion estimator is
\[
 \widehat\mu=\frac12\sum_{i=1}^2\frac{R_i y_i}{\pi_i}
             =20R_1+80R_2.
\]
Its expectation is 50 under either independent inclusion or a shared
inclusion decision $R_1=R_2$. Under independent inclusion, its four equally
likely values are $0,20,80,100$, giving variance 1700. Under shared inclusion,
it equals zero or 100 with equal probability, giving variance 2500.
The marginal probabilities and expected sample size are identical in the
two experiments~\citep{exec13_observation_exact}.

More generally, for a fixed population of size $N$, known positive
inclusion probabilities and joint probabilities
$\pi_{ij}=E(R_iR_j)$, direct expansion gives
\begin{equation}
 \operatorname{Var}(\widehat\mu)
 =\frac{1}{N^2}\sum_{i,j}
   \frac{\pi_{ij}-\pi_i\pi_j}{\pi_i\pi_j}y_i y_j,
 \qquad \pi_{ii}=\pi_i.
 \label{exec13:observation-inclusion-variance}
\end{equation}
This is a design identity, obtained by expanding the variance of a weighted
sum. It is not an estimated variance from the observed sample, and it assumes
that the origin values and selection law have the stated meaning. In the
two-origin example the extra cross term is 800. Replacing the joint law
by independent selection because a file contains only the two marginal
probabilities would discard that term.

If the second origin is never included, populations $(20,80)$ and
$(20,120)$ induce the same observed law but have means 50 and 70. The
missing information is different from an imprecise estimate of a positive
inclusion probability. No finite replacement for $1/0$ identifies the unseen
value. Outcome bounds, a changed sampling design or additional information
can restrict the answer. If inclusion probabilities are fitted, their
estimation error and relation to the outcome contribution enter a new
uncertainty argument; Equation~\eqref{exec13:observation-inclusion-variance}
does not supply that argument automatically~\citep{exec13_observation_exact}.

Identity ambiguity has another geometry. Two received records with values
20 and 80 can be mapped bijectively to two origins in either order. The
total is 100 under both maps, but the second-minus-first origin contrast is
60 or $-60$. Independently drawing a match for each record could assign
both to one origin, which is outside this stipulated bijective class.
The consumer should retain the admissible joint maps, the target functional
and any evidence used to remove a map. These two values describe an authored
measurement contrast; they are not a treatment effect or a probability
distribution over all possible linkage mechanisms.

\subsection{Noisy audit labels can change the proposed action}

Let $Z$ be a true binary adverse event and $B$ its recorded label. Write
$p=P(Z=1)$, $s=P(B=1\mid Z=1)$ and $f=P(B=1\mid Z=0)$ for prevalence,
sensitivity and false-positive probability in the declared target population.
The law of total probability gives
\begin{equation}
 q=P(B=1)=f+(s-f)p.
 \label{exec13:observation-label-law}
\end{equation}
This relation is about measurement. It does not identify a contact effect.
If sensitivity or false positives vary with contact, borrower state or time,
the corresponding conditional laws must be specified and transported
appropriately. Correcting an aggregate prevalence cannot replace the
assignment and identification assumptions needed to compare interventions.

Two strictly positive authored joint laws illustrate the information gap.
The triples $(s,f,p)=(9/10,1/10,3/10)$ and $(1/2,1/10,3/5)$ both give
$q=17/50$. More records of $B$ alone cannot distinguish these two
measurement mechanisms. An audit may resolve the issue only to the extent
that it supplies suitable independent truth, known or justified inclusion,
and a relevant population. Calling the same noisy label a gold standard
does not create that information~\citep{exec13_observation_exact}.

Suppose, solely for a planning example, that contact prevents an event with
probability one half among borrowers who would otherwise experience it,
has no adverse event effect on the others, and costs USD~200. With an event
loss of USD~1000, the assumed expected incremental value is
\[
 V(p)=500p-200\quad\text{USD per eligible borrower}.
\]
The two equally compatible reported-label laws imply values $-50$ and
100. The sign changes even before sampling error is considered. The
one-half prevention mechanism is an additional stipulated causal model;
it is not a measured treatment effect inferred from $q$. The example
therefore motivates an experiment on the consequential missing information,
without asserting measured bank returns~\citep{exec13_observation_exact}.

\subsection{Propagate the fitted measurement law and preserve singular cases}

When $d=s-f\ne0$, the algebraic inverse is $p=(q-f)/d$. Suppose a
particular sampling and fitting design justifies a joint expansion of
$(\widehat q,\widehat s,\widehat f)$ with leading contributions
$(\psi_q,\psi_s,\psi_f)$. Differentiating the inverse gives
\begin{equation}
 \psi_p=\frac{\psi_q-p\psi_s-(1-p)\psi_f}{s-f}.
 \label{exec13:observation-label-score}
\end{equation}
For a declared common normalization, write the justified covariance of
those contributions as $\Omega$. Then the corresponding leading variance
is $g^\top\Omega g$, where
$g=(1,-p,-(1-p))^\top/d$. Different audit and main-sample sizes require
their proper scaling in the joint expansion. Household dependence, shared
label adjudication, reused fitting data and selection weights cannot be
omitted merely because each marginal estimate has a reported standard error.

An exact covariance illustration makes the cross terms visible. Take
independent unit-variance $X,Y$, and stipulate
$\psi_q=X$, $\psi_s=-X+Y$, $\psi_f=0$, $p=3/10$ and $d=4/5$.
Equation~\eqref{exec13:observation-label-score} then gives leading variance
$89/32$, whereas treating the fitted measurement law as fixed gives
$25/16$. These are authored influence contributions, separate from any
claim about the sampling law of the binary example. They show why a
covariance argument must concern the whole construction. They do not
establish that arbitrary audit estimators have this expansion
\citep{exec13_observation_exact}.

When $s=f$, the recorded label probability is $f$ for every $p$.
Near this regime the inverse and its derivative can be unstable. A more
faithful proposed interface carries a region $\mathcal C$ for the joint
observable and measurement parameters and reports
\begin{equation}
 \mathcal P(\mathcal C)=
 \{p\in[0,1]:\ \exists(q,s,f)\in\mathcal C,
                    \ q=f+(s-f)p\}.
 \label{exec13:observation-label-projection}
\end{equation}
If the true triple lies in $\mathcal C$, its true prevalence belongs to
this set, by substitution. Thus a genuinely justified joint-region coverage
statement transfers to the projected set. The projection cannot manufacture
that joint guarantee from an arbitrary collection of marginal intervals.
Nor does this argument establish a sharp identified set for a richer
observed law. The observation model, candidate region and target restriction
remain part of the claim.

For fixed $s=9/10$, $f=1/10$ and an expressly deterministic sensitivity
range $q\in[3/10,19/50]$, the compatible prevalence interval is
$[1/4,7/20]$. Equivalently,
$(4p-1)(20p-7)\leq0$. Mapping that whole interval through the illustrative
value function gives $[-75,-25]$ USD. At complete loss of label information,
the physical range $p\in[0,1]$ gives $[-200,300]$ USD. An empty compatibility
set instead indicates conflict among the supplied ranges and restrictions;
it is not evidence of zero loss. The generic local quadratic service can
represent finite intervals, empty sets, the whole line and disconnected
rays, but does not impose every target's physical range. The proposed
consumer must preserve the original result and apply justified target
restrictions explicitly~\citep{exec13_observation_exact,exec13_observation_primary}.

\subsection{What the local boundary comparison can establish}

The registered comparison fixes 26 authored cases across the library and
command-line routes, with an initial call and a restarted call through each.
Two further calls retain an existing key while changing the target horizon.
All 106 slots include refused, unsupported and failed requests as well as
computed results. Complete-response parity excludes only the request
correlation identifier. This tests a bounded local contract; it does not
estimate graph accuracy, sampling coverage or bank decision value
\citep{exec13_observation_design,exec13_observation_run}.

The cases deliberately expose the distinction between a completed diagnostic
and a validated observation. An empty row array can satisfy every declared
clock vacuously. Duplicate origins, false labels and an explicitly contradicted
identity statement can accompany valid timestamps. The present clock operation
checks their timing and retains its limited qualification; it does not repair
the other defects. Asking it jointly for clocks and identity does not silently
upgrade its scope. Measurement fitting, inclusion validation and noisy-label
truth remain separate unsupported requests~\citep{exec13_observation_primary,exec13_observation_run}.

Supplied-moment calculations reveal another boundary. An unverified measurement
statement can accompany a conditional numerical answer, while an explicitly
contradicted premise causes the selected estimator to refuse. The service
does not establish the truth of either declaration, automatically require
a preceding relevant audit, or connect that audit to a full causal analysis.
The prototype recommendation is therefore a consumer admission rule that
requires a nonempty audit of the right role and version, then either verifies
the remaining premises or propagates their uncertainty and limitations.
This admission rule is specified here; it is not claimed as an implemented
bank workflow~\citep{exec13_observation_design,exec13_observation_run}.



% SOURCE docs/release/2026-10-29/source/docs/writing-repair/2026-10-29/src/cf67fcce1027578c.tex CHAR 0
All 106 registered requests matched the frozen boundary expectations, and
all 26 cases preserved semantic responses across the two routes and restart.
The denominator contains 52 computed outputs, 20 unsupported requests,
four scientific refusals, 28 failed/unavailable results and two rejected
changed-horizon requests. A computed output includes diagnostics and generic
sets; this count is not the number of identified or useful causal estimates.
All negative dispositions remain in the comparison. Independent rational
evaluation checked 1,446 points against the supplied inequalities
\citep{exec13_observation_run}.

Total invocation wall time, including failed requests, replay and worker or
CLI launch, was 42.398 seconds in this local run. Active runner
time was 43.542 seconds, including fixture import and journaling.
These measurements exclude idle gaps between bounded segments, initial
runner startup, author review and all unimplemented work. Parent and child
memory records are separate process high-water readings, not simultaneous
process-tree memory. They do not establish a bank service level, a full
workflow price or a performance advantage over another product.

The journal preserved every earlier request across four execution segments
and a final aggregation-only resume. Restart reused the saved artifact
references and idempotency identities. It did not test process termination
at every possible publication point or multi-host recovery. No native
foreign-language engine, sibling repository, live model or authentic bank
owner participated in this screen~\citep{exec13_observation_run}.


% SOURCE docs/release/2026-10-29/source/docs/writing-repair/2026-10-29/src/146d5f4fb9fee477.tex CHAR 13720


\subsection{Choose the next information purchase by its decision consequence}

The first useful prototype combines the existing clock and artifact boundary
with one explicit observation model and a consumer that preserves unresolved
maps, selection and label uncertainty. A public dictionary can explain a
recorded field, and generated truth can test whether the interface loses
information. Neither proves that a bank's historical labels are correct.
Source navigation, proposed transforms and numerical operations should remain
separately accountable when an agent or another repository supplies them.
An actual sibling trial needs its own request, source/export identity,
response, failure/restart trace and owner disposition. A contract double
cannot substitute for that execution.

The complete cost account includes source acquisition and rights, extraction,
identity resolution, audit sampling and adjudication, fitting and covariance
work, data transfer, integration, scientific review, storage, failed requests,
monitoring and recovery. Local invocation time prices only a subset of that
work. Missing prices stay unknown. Before purchasing costly bank evidence,
the proposal asks whether the current feasible decision values differ enough
that improved identity, inclusion or label information could change the action.
If all admitted values already imply the same feasible action, more precise
measurement may have little incremental decision value. When signs differ,
the relevant study targets the uncertainty causing that difference, with
its cost compared against a justified expected benefit. The present generated
examples establish this rationale for a bounded prototype and an explicit
later experiment; they do not establish that any particular bank experiment
will be profitable~\citep{exec13_observation_design,exec13_observation_exact}.
\label{exec13:observation-journey-end}


% SOURCE docs/release/2026-10-29/source/docs/writing-repair/2026-10-29/src/4a996fcba1365f41.tex CHAR 0
\section{What an executed interface result actually establishes}
\label{exec14:execution-meaning}

A proposal can contain an implemented local dispatcher, a proposed Java adapter,
a mocked external consumer and an exact mathematical reference at the same time.
The integration argument depends on keeping their evidence separate. The
scientific meaning of an output does not become stronger merely because it
passes through a real CLI or MCP server. A graph-syntax check remains a syntax
check; an unsupported policy request remains unsupported when its refusal is
correctly transported to a client.

The shared response therefore needs two independent descriptions. One records
how the result was obtained: design only, static inspection, project reference,
consumer double, actual local engine, actual transport, or an externally run
consumer. The other records the scientific disposition: computed within a
specified method, refused because declared premises conflict, unsupported,
unknown or unavailable. A generated dataset can be processed by an actual
engine. A mock can also return realistic numbers. Input provenance and execution
provenance must both be visible before a result is interpreted.

For an actual execution claim, the evidence binds the selected source and
schema, runtime and dependency versions, scientific context, input bytes,
route, request, response and terminal failure state. A run journal retains all
attempts and resource scope. An external package name or a source-code pin alone
is insufficient. Replays should be distinguished from fresh computations and
conflicting context should invalidate an idempotency key. Owner-run consumer
acceptance adds a named owner, actual consumer configuration and its acceptance
criteria; a local substitute cannot confer that acceptance.

The runtime revision summarized by this evidence implemented nine narrow operations: capabilities,
graph syntax review, supplied-DAG back-door identification, supplied-moment
Gaussian residual inversion, deterministic quadratic sensitivity sets,
observation-clock diagnostics, reports and job status/cancellation. Its actual
local transport evidence includes eleven scientific cases agreeing across
library, CLI and an MCP SDK client/server exchange using protocol 2025-11-25.
The qualification applies to the recorded local stdio environment. The separate later parity study in Section~\ref{exec16:semantic-parity} extends transport evidence for its own request set; it does not retroactively change this trial's denominator. It does not
establish modern-protocol or HTTP support, arbitrary discovery, a longitudinal
policy estimator or a live language-model backend. The local agent controller's
80 conformance checks involved zero live model calls
\citep{exec14_execution_evidence}.

A further servicing screen exercised 98 actual library and CLI invocations,
including restart/replay and changed-horizon conflicts. All 24 four-way semantic
comparisons agreed, and 964 independent rational membership checks reproduced
the returned deterministic sensitivity sets. The frozen suite nevertheless
retains a negative result: 94 of 98 expected outcomes matched. Four corrupt-input
cases returned the documented \texttt{artifact\_access} code where the authored
expectation named a different error. Each failed safely with no numerical
result or qualification. Diagnosing an expectation error does not rewrite the
original trial as a pass. The same trial shows that an actual engine can return
correctly scoped unsupported and refused results; it does not show that a
servicing estimator has been installed \citep{exec14_servicing_evidence}.

The distinction also changes cost interpretation. The servicing runner measured
35.4281 seconds summed across invocations and 36.5564 seconds of active runner
time. Parent and child peak resident memory were 45,056 and 104,448 KiB,
respectively; these are separate process maxima and cannot be added to claim a
concurrent memory peak. The timing covers this bounded local screen, not a
commercial engine, an external consumer or bank-scale deployment. A future
comparison should include setup, transfer, failures, review and recovery in the
same accounting boundary before proposing a service-level or economic benefit
\citep{exec14_servicing_evidence}.

\section{Servicing evidence, policy values and the proposed next experiment}
\label{exec14:servicing-evidence}

The servicing example has two different mathematical tasks. Exact enumeration
of a finite declared model can compare fixed policies. Estimating the value of
a policy from incomplete histories under hidden assignment information is a
harder problem. The available finite reference gives values $11/8$ for no
assistance, $2$ for assistance to everyone and $133/80$ under the stipulated
stress construction. These are generated model values, not fitted bank effects.
The library/CLI screen exercises supported graph, clock and sensitivity
operations; requests for the absent policy estimator return no policy estimate
\citep{exec14_servicing_evidence}.

The conservative precision calculation is an adverse result worth retaining.
At $n=1000$ the evaluation radius is approximately $8.242685524$, greater than
the complete single-policy range width of eight in the declared finite model.
The illustrative combined radius is approximately $9.615113642$. Such a bound
can be mathematically valid yet unhelpful for deciding between the fixed
policies. It is not evidence that collecting 1000 bank histories is a worthwhile
experiment. Better precision requires an explicitly justified sharper method
or a better information design; replacing the number with a narrow heuristic
would change the guarantee \citep{exec14_servicing_evidence}.

The missing-information examples are equally consequential. Two hidden-state
models can have identical observations but effects $.6$ and zero. Two
informative-dropout models can have identical observed histories but cohort
means $.25$ and $.75$. More observations of the same incomplete information do
not identify the missing distinction. A separate proposed continuous Gaussian
policy benchmark with quadratic reward also lies outside the bounded
finite-table theorem. It must retain its own tail and integrability analysis
rather than inherit the bounded guarantee by analogy
\citep{exec14_servicing_evidence}.

A planning scenario can still indicate what a prototype might be worth studying.
The fixed-policy contrast $2-11/8=5/8$, multiplied by 1000 customers and an
assumed USD~100 per utility unit, gives USD~62,500. Subtracting an assumed
additional full-workflow cost of USD~10,000 gives USD~52,500. Assistance costs
were already deducted inside the reward and are not subtracted again. Every
currency conversion, customer count and additional cost is stipulated. The
calculation is a scenario for sensitivity analysis, not measured bank ROI
\citep{exec14_servicing_evidence}.

The fitted servicing and policy-uncertainty studies in Chapters~\ref{gap29:servicing} and~\ref{gap29:policy-uncertainty} now supply separate generated evidence under controlled assignment, overlap and dropout mechanisms. Their findings must be read alongside the exact references and adverse bounds here; they neither install a policy estimator in this tested service revision nor resolve unrestricted hidden-state assignment. Any further policy experiment should identify the remaining mechanism or precision question that it changes. It
should keep training and evaluation separate, retain every failed request, and
compare a fixed policy with the chosen trained-policy alternative under the same
reward and observation law. Bank data become relevant only if the prototype
isolates a consequential assumption or precision requirement that the public
and generated study cannot resolve. An honest negative precision result and a typed refusal remain useful proposal evidence.

Across these journeys, the operating recommendation is to preserve the whole scientific object and every attempted request before judging convenience or speed. Local report recovery and an exact retention proof support their declared boundaries; publication races, missing dependencies, unverified covariance and unhelpful precision require distinct remedies. A reviewed manual workflow remains the meaningful comparator until an adapter demonstrates the same target, uncertainty and failure behavior at an acceptable complete cost. The following construction chapter turns to the records and transformations on which that comparison depends.

